\par\Needspace{4\baselineskip}
\begin{center}
\hypertarget{native-060-article043}{}\pdfbookmark[2]{43.}{article-043}43.
\end{center}
To illustrate the preceding investigations more fully, we add a complete example of calculating the true anomaly and radius vector from the time, using the numbers of Art. 38 again. We thus set $e=0{,}967\,6456\,7$, $\log q=9{,}765\,6500$, $t=63{,}54400$, from which we first derive the constants $\log\alpha=0{,}305\,2357$, $\log\beta=8{,}221\,7364$, $\log\gamma=0{,}002\,8755$.

\SourceLeaf{061}This gives $\log\alpha t=2{,}108\,3102$, to which the approximate value $w=99^0\,6^{\prime}$ corresponds in \textsc{Barker}'s table. From this we obtain $A=0{,}02292$ and, from our table, $\log B=0{,}000\,0040$. Hence the corrected argument with which to enter \textsc{Barker}'s table is $=\log\dfrac{\alpha t}{B}=2{,}108\,3062$, to which $w=99^0\,6^{\prime}\,13^{\prime\prime}{,}18$ corresponds. The calculation then proceeds as follows.
\begin{center}
\begin{tabular}{@{}ll@{}}
\makebox[18em][l]{$\log\operatorname{tang}\tfrac{1}{2}w^2$\hfill .\,.\,.\,.\,.\,.\,.\,.\,.} & $0{,}138\,5934$\\[.3em]
\makebox[18em][l]{$\log\beta$\hfill .\,.\,.\,.\,.\,.\,.\,.\,.\,.\,.\,.\,.\,.} & $8{,}221\,7364$\\[.3em]
\hline
\makebox[18em][l]{$\log A$\hfill .\,.\,.\,.\,.\,.\,.\,.\,.\,.\,.\,.\,.} & $8{,}360\,3298$\\[.3em]
$A$\hfill $=$ & $0{,}022\,9260\,8$\\[.3em]
\multicolumn{2}{@{}l@{}}{Hence the same $\log B$ as before}\\[.3em]
$C$\hfill $=$ & $0{,}000\,0242$\\[.3em]
$1-\tfrac{4}{5}A+C$\hfill $=$ & $0{,}981\,6833$\\[.3em]
$1+\tfrac{1}{5}A+C$\hfill $=$ & $1{,}004\,6094$
\end{tabular}
\end{center}
\begin{center}
\begin{tabular}{@{}ll@{}}
\makebox[20em][l]{$\log\operatorname{tang}\tfrac{1}{2}w$\hfill .\,.\,.\,.\,.\,.\,.\,.\,.\,.} & $0{,}069\,2967$\\[.3em]
\makebox[20em][l]{$\log\gamma$\hfill .\,.\,.\,.\,.\,.\,.\,.\,.\,.\,.\,.\,.\,.\,.} & $0{,}002\,8755$\\[.3em]
\makebox[20em][l]{$\tfrac{1}{2}\operatorname{Comp.}\log(1-\tfrac{4}{5}A+C)$\hfill .} & $0{,}004\,0143$\\[.3em]
\hline
\makebox[20em][l]{$\log\operatorname{tang}\tfrac{1}{2}v$\hfill .\,.\,.\,.\,.\,.\,.\,.\,.\,.\,.} & $0{,}076\,1865$\\[.3em]
$\tfrac{1}{2}v$\hfill $=$ & $50^0\quad0^{\prime}\quad0^{\prime\prime}$\\[.3em]
$v$\hfill $=$ & $100\quad0\quad0$\\[.3em]
\makebox[20em][l]{$\log q$\hfill .\,.\,.\,.\,.\,.\,.\,.\,.\,.\,.\,.\,.\,.\,.} & $9{,}765\,6500$\\[.3em]
\makebox[20em][l]{$2\operatorname{Comp.}\log\cos\tfrac{1}{2}v$\hfill .\,.\,.\,.\,.\,.} & $0{,}383\,8650$\\[.3em]
\makebox[20em][l]{$\log(1-\tfrac{4}{5}A+C)$\hfill .\,.\,.\,.\,.\,.\,.} & $9{,}991\,9714$\\[.3em]
\makebox[20em][l]{$\mathrm{C.}\log(1+\tfrac{1}{5}A+C)$\hfill .\,.\,.\,.\,.} & $9{,}998\,0028$\\[.3em]
\hline
\makebox[20em][l]{$\log r$\hfill .\,.\,.\,.\,.\,.\,.\,.\,.\,.\,.\,.\,.\,.\,.} & $0{,}139\,4892$
\end{tabular}
\end{center}
\noindent If the factor $B$ had been neglected altogether in this calculation, the true anomaly would have come out too large by only $1^{\prime\prime}$.
